% Propietario reproducido por unidad demostrativa; véase PROPIETARIOS_CONTINUO_UNIVERSO.json.
% Origen: XI ES CUMPLIMIENTO_EDITORIAL_20260928, sections/continuo_conjunto.tex:1--629.
% Cuerpos y pruebas conservados; sólo namespace, rutas y remisiones contextuales declaradas.
\section{Joint construction of the continuum through compatible histories and prolongations}
\label{hmtfg:base:iv:continuo-conjunto}
The arithmetic expressed by a normal form preserves a precise readout of the state that produces it. To situate that readout within HMT, the histories, sheets and operations acting on them must also be retained. This chapter brings together that common substrate. Its different constructions arise from the same tree of APP–TRIT–TPK executions; the expository divisions allow their properties to be studied without turning them into independent domains chosen by a subsequent application.

The starting point is the formal kernel presented above. An admissible emission contains the APP operation, TRIT orientation, sheet change, remainder and quotient, together with the effective TPK memory update. When only a remainder is expressed, the execution history still belongs to the starting object. Two executions with identical final readout are not identified on that basis alone.

\input{\HMTFGRoot/latex/antecedentes_comunes/supervivencia_punto_fijo.tex}

\subsection{Admissible histories and memory by composition}
Let \(H_0:=\{\ast\}\). Level \(H_1\) is the finite nonempty set of admissible APP--TRIT seeds, and, for \(k\geq1\), \(H_k\) is the set of enriched histories of length \(k\) that admit compatible TPK prolongations to every finite horizon. Its elements are constructed by successive emissions from the initial domain of the kernel. We write \(\zeta\varepsilon\) for the prolongation of \(\zeta\) by an admissible emission \(\varepsilon\), and \(\rho_k(\zeta\varepsilon)=\zeta\) for deletion of the last emission. The label produced by the emission is retained, even when different labels express the same scalar state. The space of complete histories is defined by
\[
 H_\infty^{\rm enr}:=\varprojlim_k(H_k,\rho_k).
\]
The cofinal subsystem \(X_N:=H_{6N}\), with \(X_0=H_0\), agrees with the sextet clock used in the cardinality proof; cofinality gives a canonical identification \(H_\infty^{\rm enr}\cong\varprojlim_N(X_N,\rho_{6(N+1),6N})\).

On the two arithmetic sheets, the local data satisfy the exact divisions

\[
\delta_\diamond=r_\diamond+9q_\diamond,
\qquad 0\leq r_\diamond<9,
\qquad \diamond\in\{\Sigma,\Pi\}.
\]
The oriented TRIT coordinate likewise admits an expression \(\Delta=o_r+3\kappa\), with the oriented representative declared by the kernel. Preserving the quotient allows those representations to be composed without replacing an emission by its isolated remainder.

In a memory coordinatization, the update of an emission is written

\[
U_\varepsilon(m)=C_\varepsilon m+\kappa_\varepsilon.
\]
For a trajectory \(\alpha=(\varepsilon_1,\ldots,\varepsilon_k)\), the effective operation is the ordered composition of these updates. Its linear part and its translation are

\[
C_\alpha=C_{\varepsilon_k}\cdots C_{\varepsilon_1},
\qquad
\kappa_\alpha=
\sum_{j=1}^{k}C_{\varepsilon_k}\cdots
C_{\varepsilon_{j+1}}\kappa_{\varepsilon_j}.
\]
The empty product in the last sum is the identity.

\begin{proposition}[Affine compatibility of concatenation]
\label{hmtfg:base:prop:continuo-compatibilidad-afin}
If \(\alpha\) precedes \(\beta\), then

\[
C_{\beta\alpha}=C_\beta C_\alpha,
\qquad
\kappa_{\beta\alpha}=C_\beta\kappa_\alpha+\kappa_\beta,
\qquad
U_{\beta\alpha}=U_\beta U_\alpha.
\]
\end{proposition}
\begin{proof}
Substituting \(U_\alpha(m)\) into \(U_\beta\) produces \(C_\beta C_\alpha m+C_\beta\kappa_\alpha+\kappa_\beta\). Decomposition into linear part and translation gives the three identities. Associativity of composition proves that the result does not depend on an auxiliary subdivision of the trajectory.
\end{proof}

This cocycle identity is the precise content of accumulated memory. The quantity transported by a second trajectory depends on the coordinatization update it performs; in general, it is not the untransported sum of two records.

The nonadic return preserves a second distinction. In phase and turn coordinates, with \(\bar r\in\{0,\ldots,8\}\), its representation is

\[
\sigma_9(w,r)=
\left(w+\left\lfloor\frac{\bar r+1}{9}\right\rfloor,
r+1\pmod9\right),
\qquad
\sigma_9^9(w,r)=(w+1,r).
\]
During nine steps, exactly one crossing from \(8\) to \(0\) occurs. The phase returns and the turn coordinate increases. The enriched holonomy also retains the memory updates accompanying that path.

\subsection{Tree refinement and measure of histories}
The initial measure is the canonical census measure: \(\mu(\ast)=1\) and \(\mu(c)=1/|H_1|\) for each seed \(c\in H_1\). In this section each history has a finite positive number of prolongations. Deletion of emissions defines the refinement maps. Admissibility of the kernel supplies the prolongations; when a realization uses a neutral emission, this explicitly gives a section of \(\rho_k\). Existence of that section refers to the tree containing that emission, not to every arbitrary restriction of routes.

Let \(d(\zeta)\geq1\) be the number of admissible emissions at \(\zeta\). The APP order and the census of prolongations determine the canonical law \(p_{\rm can}(\varepsilon\mid\zeta)=1/d(\zeta)\). Without additional parameters, this law defines

\[
\mu(\zeta\varepsilon)=\frac{\mu(\zeta)}{d(\zeta)}.
\]
A nonuniform weighting may replace \(p_{\rm can}\) only when it has already been produced and added to the enriched record. The correlative constructor of this chapter uses only \(p_{\rm can}\); its measure is therefore a function of the APP--TRIT--TPK tree and not of an external choice.

\begin{proposition}[Cylindrical compatibility and conditional expectation]
\label{hmtfg:base:prop:continuo-esperanza-condicional}
At each level we have

\[
\sum_{\rho_k(y)=x}\mu(y)=\mu(x).
\]
In the spaces \(L^2(H_k,\mu_k)\), the maps

\[
(\mathsf J_kf)(y)=f(\rho_k(y)),
\qquad
(E_kg)(x)=\sum_{\rho_k(y)=x}\frac{\mu(y)}{\mu(x)}g(y)
\]
satisfy \(E_k=\mathsf J_k^*\), \(E_k\mathsf J_k=I\) and \(\|\mathsf J_kf\|=\|f\|\).
\end{proposition}
\begin{proof}
The first identity is the normalization of local probabilities. Applying it to \(|f(x)|^2\) gives the isometry. Grouping the sum by fibers of \(\rho_k\) gives the formula for \(E_k\); on a function constant on each fiber, that formula returns its value.
\end{proof}

The compatible masses define a probability measure on the space of infinite histories: the cylinder of a finite history receives its mass \(\mu(\zeta)\). The measure extends from the cylinder algebra, and its uniqueness follows because cylinders generate the \(\sigma\)-algebra. If the tree is finitely branching and every history has a prolongation, the inverse limit of its levels under emission deletion is nonempty. In the cylinder topology, every nonempty cylinder has positive measure.

This measure distinguishes histories that converge to the same representation. The uniform measure on terminal values would be a different readout, since it would remove the multiplicities that the construction has just preserved.

\input{\HMTFGRoot/latex/antecedentes_comunes/ch_memoria_y_naturalidad.tex}

\subsection{Transport by sheets and exact accounting of outputs}
First consider the space with an orthonormal basis formed by the histories at level \(k\). The transport preserving the label of each prolongation is

\[
\mathcal U_ke_\zeta=
\sum_{\varepsilon\ \mathrm{admisible\ en}\ \zeta}
\sqrt{p_{\rm can}(\varepsilon\mid\zeta)}\,e_{\zeta\varepsilon}.
\]
The successor fibers of distinct antecedents are disjoint. Normalization of \(p_{\rm can}\) therefore proves \(\mathcal U_k^*\mathcal U_k=I\). The representation by history masses is related to this one through the diagonal change of basis incorporating \(\sqrt{\mu(\zeta)}\).

The sheet label is part of the complete register of a history. We write
\[
\widehat\zeta=((\zeta_0,s_0),\ldots,(\zeta_k,s_k)),\qquad
s_j\in\{\Sigma,\Pi\},
\]
where \(s_j\) is the active sheet at step \(j\). Prolongation adds the emission and the sheet \(s_{k+1}\) determined by the operational TPK selector, preserving the entire prefix, including the initial sheet. In the basis of these lifted histories,
\[
\mathcal U_ke_{\widehat\zeta}
=\sum_{\varepsilon\ \mathrm{admisible\ en}\ \widehat\zeta}
\sqrt{p_{\rm can}(\varepsilon\mid\widehat\zeta)}\,e_{\widehat\zeta\varepsilon},
\qquad
P_ke_{\widehat\zeta}=\mathbf1_{s_k=\Sigma}e_{\widehat\zeta},\quad
Q_k=I-P_k.
\]
Thus \(P_k\) reads the current sheet, while histories remain distinguished by their complete prefix. Two initial histories reaching the same active sheet retain distinct successor fibers. The isometry proof through disjoint collections of prolongations applies literally to this basis.
Tritic orientation is another coordinate of the record: the value \(0\) preserves the threshold readout on each sheet and does not add a third sheet to this decomposition. We now separate sheet preservation and sheet exchange:

\[
A_k=P_{k+1}\mathcal U_kP_k+Q_{k+1}\mathcal U_kQ_k,
\qquad
\eta_k=Q_{k+1}\mathcal U_kP_k+P_{k+1}\mathcal U_kQ_k.
\]
\begin{proposition}[Gram identity by sheets]
\label{hmtfg:base:prop:continuo-gram-hojas}
For a linear transport \(U\), whether isometric or not, the preceding decomposition satisfies

\[
A^*A+\eta^*\eta=PU^*UP+QU^*UQ.
\]
In particular, for \(U=\mathcal U_k\),

\[
A_k^*A_k+\eta_k^*\eta_k=I.
\]
\end{proposition}
\begin{proof}
In \(A^*A\) and \(\eta^*\eta\), mixed terms containing \(P_{k+1}Q_{k+1}\) vanish. Summing the remaining terms yields \(P_kU^*(P_{k+1}+Q_{k+1})UP_k\), and its analogue with \(Q_k\). The first identity follows from complementarity of the projectors; the second also uses \(U^*U=I\).
\end{proof}

The first identity is the general formula. Replacing its right-hand side by \(U^*U\) requires the cross blocks of that Gram matrix to vanish. Isometric preservation of history transport provides that condition here; it should not automatically be attributed to a subsequent compression.

Write \(F_{0,0}=I\), \(F_{0,n}=A_{n-1}\cdots A_0\) and \(T_N=F_{0,N}^*F_{0,N}\).

\begin{theorem}[Finite decomposition with terminal residual]
\label{hmtfg:base:thm:continuo-descomposicion-terminal}
For each horizon \(N\),

\[
I=\sum_{n=0}^{N-1}
F_{0,n}^*\eta_n^*\eta_nF_{0,n}+T_N.
\]
The sequence \(T_N\) is positive and decreasing, and has a positive strong limit \(T_\infty\). The sum of outputs converges strongly to \(I-T_\infty\).
\end{theorem}
\begin{proof}
Proposition~\ref{hmtfg:base:prop:continuo-gram-hojas} gives \(T_n-T_{n+1}=F_{0,n}^*\eta_n^*\eta_nF_{0,n}\). The telescoping sum proves the finite equality and monotonicity. For each vector, the quadratic forms decrease to a limit; polarization defines a bounded positive form and, by the representation theorem for bounded forms, an operator \(T_\infty\). Convergence is strong because, for a positive operator \(S\leq I\), \(\|Sv\|^2\leq\langle v,Sv\rangle\); this is applied to \(S=T_N-T_\infty\).
\end{proof}

The terminal residue is retained until its extinction in the domain under consideration is proved. This distinction prevents a phase return, by itself, from becoming an assumption of exhaustion of all histories.

\subsection{Oriented balance and cancellation under refinement}
A history carries its accumulated orientation \(o(h)\) and the counts of its visits to the two sheets. The corresponding local balance is

\[
b(h)=o(h)\bigl(N_\Sigma(h)-N_\Pi(h)\bigr),
\qquad
b^+=\frac{|b|+b}{2},\quad b^-=\frac{|b|-b}{2}.
\]
When a reader expresses the mean over prolongations, define \(b_c=Eb_f\). Orientation remains in the argument of \(b_f\); it is not removed before averaging.

\begin{proposition}[Exact cancellation defect]
\label{hmtfg:base:prop:continuo-defecto-cancelacion}
The quantity

\[
\kappa_{\mathrm{can}}=
\frac{E|b_f|-|Eb_f|}{2}
\]
is nonnegative and satisfies

\[
E(b_f^+)=(b_c)^++\kappa_{\mathrm{can}},
\qquad
E(b_f^-)=(b_c)^-+\kappa_{\mathrm{can}}.
\]
\end{proposition}
\begin{proof}
The weighted triangle inequality gives \(|Eb_f|\leq E|b_f|\). Writing the positive and negative parts through \((|b|\pm b)/2\) yields both equalities.
\end{proof}

The information omitted when only the balance is expressed is the common cancellation of opposite contributions. This quantity is determined from the same prolongations used by the preceding transport. The factor \(1/9\) appears only when the fiber contains nine prolongations with equal weights; conditional expectation is the rule that remains valid for general branching and weightings.

\subsection{Ternary incidence of the sheets and rectangular complex}
The two sheets determine the module \(V=\mathbb F_3e_\Sigma\oplus\mathbb F_3e_\Pi\). Its four projective directions are the lines generated by \(e_\Sigma,e_\Pi,e_\Sigma+e_\Pi,e_\Sigma-e_\Pi\). They give rise to six unordered pairs. The eight nonzero vectors are their oriented representatives. Thus, the counts \(4,6,8\) belong to distinct operations on the same sheet module.

On the six pair positions, we define

\[
(Ba)_{\{L,L'\}}=a_L+a_{L'},
\qquad
s(q)=\sum_{\{L,L'\}}q_{\{L,L'\}}.
\]
Each direction occurs in three pairs; hence \(sB=0\) in \(\mathbb F_3\). The reader \(W(q)=\mathbf1_{s(q)=0}-1/3\) is invariant under \(q\mapsto q+Ba\). On the uniform ambient space \(\mathbb F_3^6\), its moments are

\[
\mathbb EW=0,\qquad
\mathbb EW^2=\frac29,\qquad
\mathbb EW^3=\frac2{27}.
\]
Indeed, \(s\) is surjective and each fiber contains \(3^5\) elements. The reader equals \(2/3\) with probability \(1/3\) and \(-1/3\) with probability \(2/3\). These means describe the uniform ambient space; the distribution induced by a subset of histories must be calculated with its own multiplicities, already preserved in \(H_k\).

To fix the ports without losing multiplicities, at depth \(k\) retain the labeled readouts \(\lambda_\Sigma(\zeta)=(\zeta,\Sigma)\) and \(\lambda_\Pi(\zeta)=(\zeta,\Pi)\). Their sets are
\[
\mathscr P_k^\Sigma=\{\lambda_\Sigma(\zeta):\zeta\in H_k\},\qquad
\mathscr P_k^\Pi=\{\lambda_\Pi(\zeta):\zeta\in H_k\}.
\]
They are finite and nonempty because \(H_k\) is. Values, quotients and remainders are expressed on those ports; equality of values does not identify their labels. The product \(\mathscr P_k^\Sigma\times\mathscr P_k^\Pi\) is the domain of the rectangular readout, within which the support of admissible incidences is recorded.

With this level fixed, write \(I=\mathscr P_k^\Sigma\), \(J=\mathscr P_k^\Pi\) and \(A_m=\mathbb Z/3^m\mathbb Z\). The APP lexicographic order determines the minimal surviving branch and, on it, the reference ports \(i_0,j_0\). No external choice is added: changing coordinates produces a canonically isomorphic complex. The maps

\[
\kappa(u,v)_{ij}=u_i-v_j,
\qquad
(\delta w)_{ij}=w_{ij}-w_{ij_0}-w_{i_0j}+w_{i_0j_0}
\]
define the following sequence, valid over the finite ring \(A_m\):

\[
0\longrightarrow A_m\longrightarrow
A_m^I\oplus A_m^J\mathrel{\mathop{\longrightarrow}^{\kappa}}
A_m^{I\times J}\mathrel{\mathop{\longrightarrow}^{\delta}}
A_m^{(I\setminus\{i_0\})\times(J\setminus\{j_0\})}
\longrightarrow0.
\]
\begin{proposition}[Exactness of rectangular incidence]
\label{hmtfg:base:prop:continuo-exactitud-rectangular}
The preceding sequence is exact, with first arrow given by the constant diagonal.
\end{proposition}
\begin{proof}
The equality \(u_i-v_j=0\) for each pair forces all coordinates of \(u\) and \(v\) to be the same constant. Direct substitution gives \(\delta\kappa=0\). If \(\delta w=0\), take \(u_i=w_{ij_0}\) and \(v_j=w_{i_0j_0}-w_{i_0j}\). Then \(u_i-v_j=w_{ij}\). Finally, any matrix of the last term is extended with zero reference row and column; its image under \(\delta\) is the initial matrix. The proof uses only addition and subtraction, so it does not require division by a unit of the ring.
\end{proof}

Reducing from \(A_{m+1}\) to \(A_m\) while keeping \(k,\mathscr P_k^\Sigma,\mathscr P_k^\Pi,i_0,j_0\) fixed commutes with both maps. The reconstruction formulas are the same at each modular depth. This is the arithmetic compatibility proved in this section. The prolongation of histories \(k\to k+1\) preserves prefix labels; its maps between port sets are a datum distinct from the ring change \(m\to m+1\).

These maps are obtained from emission removal:
\[
\rho_k^\Sigma(\zeta\varepsilon,\Sigma)=(\zeta,\Sigma),\qquad
\rho_k^\Pi(\zeta\varepsilon,\Pi)=(\zeta,\Pi).
\]
The minimal surviving branch of the APP order just constructed has compatible prolongations. Its two labels define compatible reference ports at every level. Coordinate transport is precomposition:
\[
(L_ku)_{i'}=u_{\rho_k^\Sigma(i')},\quad
(L_kv)_{j'}=v_{\rho_k^\Pi(j')},\quad
(L_kw)_{i'j'}=w_{\rho_k^\Sigma(i'),\rho_k^\Pi(j')}.
\]
For the last module of the sequence, first extend the matrix by zero on the reference row and column, apply this same formula and restrict to the new complement of the reference ports. This rule defines \(L_k\) even when a new port has the marked port as its antecedent.

\begin{proposition}[Naturality of the rectangular complex under prolongation]
\label{hmtfg:base:iv:rectangular-natural}
The preceding maps satisfy
\(L_k\kappa_k=\kappa_{k+1}L_k\) and
\(L_k\delta_k=\delta_{k+1}L_k\).
They also commute with coefficient reduction
\(A_{m+1}\to A_m\) and compose according to prefix removal.
\end{proposition}
\begin{proof}
The first identity follows from \(u_{\rho_k^\Sigma(i')}-v_{\rho_k^\Pi(j')}\). In the second, each of the four terms defining \(\delta\) is transported by the same precomposition. Compatibility of the marked ports identifies their rows and columns; when an antecedent is the marked port, the four terms cancel and agree with extension by zero. All formulas are sums and differences with integer coefficients, so they commute with modular reduction. Composition of precompositions is precomposition with the composite prefix, which proves the last assertion.
\end{proof}

\subsection{Memory complex and compatibility of incidences}
The same accumulated memory coordinate \(c_y\in\mathbb Z\), in a residence \(y\), defines a free integer complex. Its reduction to \(A_m\) uses the same integer register before reducing it; successive depths are therefore compatible. Its generators are \(f_y\) in degree \(2\), \(e_y,b_y,r_y\) in degree \(1\), and \(g_y\) in degree \(0\):

\[
\partial f_y=3c_y e_y+b_y,\qquad
\partial e_y=g_y,\qquad
\partial b_y=-3c_y g_y,\qquad
\partial r_y=0.
\]
The identity \(\partial^2=0\) is verified on \(f_y\) through \(3c_yg_y-3c_yg_y=0\), and is immediate on the other generators.

If a trajectory \(\alpha:y\to y'\) updates memory, its transport in the complex sends the correspondingly named generators to their new representatives, except

\[
\Psi_\alpha(b_y)=b_{y'}+3(c_{y'}-c_y)e_{y'}.
\]
\begin{proposition}[Naturality of the memory complex]
\label{hmtfg:base:prop:continuo-naturalidad-memoria}
We have \(\partial\Psi_\alpha=\Psi_\alpha\partial\) and \(\Psi_{\beta\alpha}=\Psi_\beta\Psi_\alpha\). The maps extend linearly over the chosen coefficient ring.
\end{proposition}
\begin{proof}
On \(f_y\), the image of its boundary is \(3c_ye_{y'}+b_{y'}+3(c_{y'}-c_y)e_{y'} =3c_{y'}e_{y'}+b_{y'}\). On \(b_y\), the boundary of its image is \(-3c_{y'}g_{y'}+3(c_{y'}-c_y)g_{y'}=-3c_yg_{y'}\). The other two cases are immediate. In a concatenation, memory differences sum telescopically, proving the second identity.
\end{proof}

The relation between two representatives of the same residence can also be written as an explicit boundary. For an integer coordinate \(v\), reduced in the same ring where appropriate, let

\[
z^-=r_y,\qquad z^+=r_y+3c_yv e_y+v b_y,
\qquad h_*=-vf_y.
\]
Then \(\partial h_*=z^--z^+\). Under a transport that also updates \(v\), the correction \(K_\alpha=(v_{y'}-v_y)f_{y'}\) satisfies \(K_{\beta\alpha}=K_\beta+\Psi_\beta K_\alpha\). The symbols \(K_\alpha\) of this local homotopy do not denote the dodecaphasic seal \(K\): they belong to another domain and their function is specified by the preceding identity. Their function can also be checked in the equality

\[
z^+_{y'}-\Psi_\alpha z^+_y=\partial K_\alpha.
\]
Both sides equal \((v_{y'}-v_y)(3c_{y'}e_{y'}+b_{y'})\).

\begin{proposition}[Homology and memory of the local complex]
\label{hmtfg:base:prop:continuo-homologia-memoria}
Over \(\mathbb Z\), or after reduction to \(A_m\), we have

\[
H_0(\Gamma)=H_2(\Gamma)=0,
\qquad H_1(\Gamma)\cong A_m[r_y]
\]
in the modular realization, and \(H_1(\Gamma)\cong\mathbb Z[r_y]\) in the integer realization.
\end{proposition}
\Needspace{9\baselineskip}
\begin{proof}
Define the homotopy of degree \(1\) by \(h(g_y)=e_y\), \(h(b_y)=f_y\), and zero on the other generators. Let \(p\) be the projection onto the summand generated by \(r_y\) and \(\iota\) its inclusion. Checking on each generator gives

\[
\partial h+h\partial=I-\iota p.
\]
In particular, on \(b_y\) the terms \(3c_ye_y\) and \(-3c_ye_y\) cancel. The complex deformation retracts onto the module generated by \(r_y\), concentrated in degree \(1\), which proves the assertion.
\end{proof}

The transport \(\Psi_\alpha\) is a shear of determinant \(1\). Thus, this complex explicitly preserves how \(c_y\) changes at the chain level, but its local homology does not distinguish the values of that coordinate. The full memory remains in the labels and transports of the source state; it is not recovered from \(H_1\) alone. A memory-dependent numerical torsion requires the specified assembly and determinantal reading, in addition to this local complex.

\subsection{Assembly with terminal paths}
For a finite horizon, the residences and their prolongations form a path category. Each residence \(\nu\) is assigned the free module \(W(\nu)=A_m[\operatorname{Path}(\nu,\mathrm{Term})]\), which preserves each terminal continuation as a generator. An initial path acts on \(W\) by precomposition. The complex \(\Gamma(\nu)\) of the preceding section acts covariantly through \(\Psi\).

The assembly uses both transports in the same bigraded module:

\[
\mathcal B_{q,r}=
\bigoplus_{\nu_0\to\cdots\to\nu_q}
W(\nu_q)\otimes_{A_m}\Gamma_r(\nu_0).
\]
The normalized path complex is used: degenerate chains containing identities are omitted. The finite horizon then bounds the number of strict prolongations and the horizontal degree. The horizontal boundary removes an intermediate residence by composing the adjacent arrows. At the first endpoint it transports \(\Gamma\); at the last it precomposes \(W\). The alternating signs yield the path boundary \(d_{\mathrm{bar}}\). The total differential is

\[
D=d_{\mathrm{bar}}+(-1)^q\partial_\Gamma.
\]
The face identities cancel the terms of \(d_{\mathrm{bar}}^2\) in pairs. The naturality of Proposition~\ref{hmtfg:base:prop:continuo-naturalidad-memoria} implies that \(d_{\mathrm{bar}}\) commutes with \(\partial_\Gamma\); the factor \((-1)^q\) cancels the cross terms of the square. Therefore, \(D^2=0\). Compatibility is proved on path generators and is not introduced by a label of membership in a complex.

Terminal evaluation likewise does not require choosing a particular history. For each continuation \(h\), the effective composition of its updates is applied. The result preserves the pair formed by \(h\) and its reading, with its multiplicity. If \(k<\ell<N\), the continuations decompose as a disjoint union of prefixes up to \(\ell\) and their continuations up to \(N\). Proposition~\ref{hmtfg:base:prop:continuo-compatibilidad-afin} proves that evaluating directly or evaluating these two parts produces the same reading. This yields the naturality square of the terminal.

The path modules, incidences, memory transport and their differentials arise from the same prolongations. Their assembly does not consist in adding to a state five coordinates that would be preserved by definition: their relations are constructed through sums, boundaries and effective compositions, and are verified on their generators.

\subsection{Determinants, limit and arithmetic reading}
In each realization, the local free complex \(\Gamma\) admits its graded determinant line

\[
\operatorname{Det}\Gamma=
\bigotimes_r\left(\bigwedge^{\mathrm{rk}\Gamma_r}\Gamma_r\right)^{(-1)^r}.
\]
The total assembly of the preceding section, in turn, has the modules \(\operatorname{Tot}_d\mathcal B=\bigoplus_{q+r=d}\mathcal B_{q,r}\). At a finite horizon and over the same ring, they are free of finite rank and only finitely many of them are nonzero. Its line is
\[
\operatorname{Det}(\operatorname{Tot}\mathcal B)=
\bigotimes_{q,r}
\left(\bigwedge^{\operatorname{rk}\mathcal B_{q,r}}
\mathcal B_{q,r}\right)^{(-1)^{q+r}}.
\]
This formula follows from decomposition as a direct sum in each degree. It distinguishes the line of the local complex from that of the assembly with all paths; their bases are ordered through the preserved labels. Modular reduction commutes with these exterior powers because the modules and their bases come from the same free integer modules at a fixed horizon.
The line is defined by the modules and their degrees, both over \(\mathbb Z\) and after reduction to \(A_m\); the negative exponent denotes the dual module. A nonzero numerical torsion additionally requires specifying the complex over a field, the bases and the relevant homology identifications. To calculate Smith factors, the integer matrix produced before reduction is used, not an arbitrary lift of a modular matrix. A square integer block with nonzero determinant is invertible over \(\mathbb Q\); its reduction is invertible over \(A_m\) only when that determinant is a unit, that is, when it is not divisible by \(3\). Refinement of memories and incidences allows the stability of the relevant invariant factors to be formulated and checked; these conditions do not follow from the existence of the determinant line.

Passage to the limit preserves the compatible projective systems of remainders, histories and their multiplicities. The squares proved in the preceding sections remain commutative because their maps are given by the same formulas at each level. For an Archimedean coordinate, the additional step consists in expressing nonempty cylinder intervals, compatible by inclusion and with diameters tending to zero. Their nested compact closures have a one-point intersection: existence follows by compactness and uniqueness by the diameter bound. That is the value of the limit readout. If the intervals are half-open, the point may belong only to their closures; the boundary representation rule then determines its positional expression. The value does not replace the history that produced it. The other representation preserves the incidences of the same state. Both destinations share their origin without thereby receiving the same information.

% Procedencia: continuo_conjunto.tex:434--444; remisiones contextuales a
% publicaciones posteriores, no premisas de las pruebas del continuo ni de Feigenbaum.
% Se conserva la descripción positiva sin exigir etiquetas de un anfitrión concreto.
The coinductive regions, the dodecaphase record and exceptional incidence are expressed from that substrate through their respective maps. The Hilbert-space realization subsequently preserves phase labels and specifies its own unital inclusion; duality and screens act on the codomains constructed there. The joint operations brought together here transmit histories, sheets, orientation and memory to those representations. Identifying a particular functional with an external analytic form in turn requires writing its action on the generators and verifying its inner products: the conservation identities of this chapter do not replace that composition.

\subsection{Mathematical scope of the joint construction}
The state, tree and measure constructions; the correlative operations; and the terminal assembly with naturality and limit are developed in full in this article. In this chapter their balance, incidence and memory identities have been formulated and proved on the same generators. The five constructions thus remain operations of a single tree. Each specialization keeps its exact hypotheses visible: transport isometry, canonical weight law, treatment of the terminal remainder and nonsingularity of the determinantal block where applicable.

The preceding equations prove the memory cocycles, weighted cancellation, rectangular exactness and compatibility of the memory complex at every finite horizon, as well as their compatible passage to the limit. A subsequent analytic identification must compose its operator on these generators and verify the additional properties of its codomain; it neither alters nor replaces the joint construction proved here.

\iffalse
% Formulación preliminar conservada fuera del ensamblado; sustituida por la
% versión tipada y verificada que se carga al final de este bloque.
\subsection{Joint emergence of the five constructions}
\label{hmtfg:base:sec:cinco-construcciones-continuo-ch}

The preceding operations are now brought together without turning them into five branches. For a surviving history \(h\in X_k\), a horizon \(N>k\) and a modular depth \(m\geq1\), let the common base record be
\[
 \mathcal B_{k,N}(h):=
 \bigl(\mathscr T_{k,N}(h),
       (\varepsilon_\alpha,s_\alpha,o_\alpha,
        r_{\Sigma,\alpha},q_{\Sigma,\alpha},
        r_{\Pi,\alpha},q_{\Pi,\alpha},
        \operatorname{Carr}_\alpha,
        \kappa_\alpha,\partial_\alpha,
        \operatorname{Surv}_\alpha)_\alpha;
       \operatorname{src},\operatorname{tgt},\circ,\rho\bigr),
 \tag{C1}
 \label{hmtfg:base:eq:base-comun-registros}
\]
where \(\mathscr T_{k,N}(h)\) is the subtree of prolongations of \(h\) up to horizon \(N\). Each component of the record is an output of the APP divisions, TRIT orientation or TPK updates described in the preceding sections; none originates in a classical problem chosen afterward.

On this same base, define the five structured objects
\[
\begin{aligned}
 \mathcal F_{\rm sp}(\mathcal B_{k,N})
   &:=(\ell^2(X_j,\mu_j),\mathcal U_j,A_j,\eta_j,T_j)_{k\leq j\leq N},\\
 \mathcal F_{\rm sol}(\mathcal B_{k,N})
   &:=(b,\kappa_{\rm can},C_\alpha,\kappa_\alpha)_{\alpha\subset\mathscr T_{k,N}},\\
 \mathcal F_{\rm gau}(\mathcal B_{k,N})
   &:=(I_j,J_j,B_j,\kappa_j,\delta_j)_{k\leq j\leq N},\\
 \mathcal F_{\rm coh}^{(m)}(\mathcal B_{k,N})
   &:=(\Gamma_y\otimes A_m,\partial,\Psi_\alpha,H_\bullet)_{y,\alpha\subset\mathscr T_{k,N}},\\
 \mathcal F_{\rm det}^{(m)}(\mathcal B_{k,N})
   &:=(\operatorname{Tot}\mathcal B,\operatorname{Det}(\operatorname{Tot}\mathcal B))_{k,N,m}.
\end{aligned}
 \tag{C2}
\]
The discriminants \({\rm sp},{\rm sol},{\rm gau},{\rm coh},{\rm det}\) name, respectively, the spectral, solenoidal, gauge, cohomological and determinantal readouts of a single object. They are not input domains. If \(p_T\) forgets the structure added by reader \(T\) and returns the record \(\mathcal B_{k,N}(h)\), write \(\widetilde{\mathcal F}_T=(\mathcal B_{k,N},\mathcal F_T,p_T)\). The correlative output is the fiber product
\[
\begin{aligned}
 \mathcal J_{k,N}^{(m)}(h):={}&
 \widetilde{\mathcal F}_{\rm sp}
 \mathop{\times}_{\mathcal B_{k,N}(h)}
 \widetilde{\mathcal F}_{\rm sol}
 \mathop{\times}_{\mathcal B_{k,N}(h)}
 \widetilde{\mathcal F}_{\rm gau}\\
 &\mathop{\times}_{\mathcal B_{k,N}(h)}
 \widetilde{\mathcal F}_{\rm coh}^{(m)}
 \mathop{\times}_{\mathcal B_{k,N}(h)}
 \widetilde{\mathcal F}_{\rm det}^{(m)}.
\end{aligned}
 \tag{C3}
 \label{hmtfg:base:eq:constructor-correlativo}
\]
Equality of the five projections onto \(\mathcal B_{k,N}(h)\) requires preservation of the same tree, the same edges and the same record; this is why \eqref{hmtfg:base:eq:constructor-correlativo} is not a Cartesian product of five retrospectively assembled objects.

Deletion of the last layer of the tree induces \(\operatorname{Res}_{N+1,N}\). On the generators, the preceding propositions give the squares
\[
\begin{gathered}
 \operatorname{Res}_{N+1,N}\mathcal F_T(\mathcal B_{k,N+1})
 =\mathcal F_T(\mathcal B_{k,N}),
 \qquad T\in\{{\rm sp},{\rm sol},{\rm gau},{\rm coh}\},\\
 \operatorname{Det}(\operatorname{Tot}\mathcal B_{k,N+1})
 \simeq
 \operatorname{Det}(\operatorname{Tot}\mathcal K_{N+1})
 \otimes
 \operatorname{Det}(\operatorname{Tot}\mathcal B_{k,N}),\\
 (\partial\Psi_\alpha=\Psi_\alpha\partial),\qquad
 (L_k\kappa_k=\kappa_{k+1}L_k),\qquad
 (L_k\delta_k=\delta_{k+1}L_k).
\end{gathered}
 \tag{C4}
 \label{hmtfg:base:eq:naturalidad-D}
\]
The change \(A_{m+1}\to A_m\) also commutes with \(\partial\), \(\Psi\), \(\kappa\), \(\delta\) and the determinant lines. A single pro-system in \((N,m)\) is therefore obtained.

Let \(\mathcal X_\chi\) be the limit of compatible enriched histories. For \(\xi\in\mathcal X_\chi\), denote by \(\mathfrak L_N(\xi)\) the ordered prefix of its \(N\) complete emissions. The joint horizon fibers are
\[
 \mathfrak J_{N,\bullet}^{(m)}
 :=\coprod_{h\in X_0}
 \{(\mathfrak L,J):\mathfrak L\text{ is a branch of length }N,
       \ J=\mathcal J_{0,N}^{(m)}(h)\}.
 \tag{C5}
 \label{hmtfg:base:eq:fibras-conjuntas-horizonte}
\]
The mark \(\mathfrak L\) does not retain a tautological copy of \(\xi\): it is formed by the emissions and updates produced at each edge. The joint generator is
\[
\begin{aligned}
 \operatorname{Gen}_{\rm cont}:\mathcal X_\chi
   &\longrightarrow\varprojlim_{N,m}\mathfrak J_{N,\bullet}^{(m)},\\
 \xi&\longmapsto
   (\mathfrak L_N(\xi),\mathcal J_{0,N}^{(m)}(h_0))_{N,m},\\
 \mathcal C_{\rm cont}^{\rm disc}
   &:=\operatorname{im}(\operatorname{Gen}_{\rm cont}).
\end{aligned}
 \tag{C6}
 \label{hmtfg:base:eq:continuo-discreto}
\]

\begin{theorem}[Correlative emergence of the five constructions]
\label{hmtfg:base:thm:correlacion-cinco}
For every surviving history \(h\in X_k\), every horizon \(N\geq k+1\) and every depth \(m\geq1\), the fiber product \eqref{hmtfg:base:eq:constructor-correlativo} exists and is nonempty. Its five components emerge from the same APP--TRIT--TPK record, satisfy the naturality squares \eqref{hmtfg:base:eq:naturalidad-D}, base changes and passage to the limit. The image \(\mathcal C_{\rm cont}^{\rm disc}\) contains five inseparable aspects of the same continuum, not five continua or five independent maps.
\end{theorem}

\begin{proof}
Prolongability of the tree provides at least one edge at each interior vertex. The record of that edge determines isometric transport and its sheet decomposition, the oriented cocycle, the incidence ports, the memory complex and its determinant line. Thus \(\mathcal F_{\rm sp}\), \(\mathcal F_{\rm sol}\), \(\mathcal F_{\rm gau}\), \(\mathcal F_{\rm coh}^{(m)}\) and \(\mathcal F_{\rm det}^{(m)}\) are successively constructed on the same base, proving nonemptiness of \eqref{hmtfg:base:eq:constructor-correlativo}. None of the five constructions appears as an input coordinate.

The affine compatibility proposition proves memory composition; the Gram identity, rectangular exactness and naturality of the complex prove the first four lines of \eqref{hmtfg:base:eq:naturalidad-D}. Decomposition of the bicomplex structure by the last layer proves the multiplicative identity of the determinant line. All these operations are computed on the same generators and commute with modular reduction. Fiber, relation, number, orientation, carry, memory, boundary and route remain in \(\mathfrak L_N(\xi)\); multisection, terminals, and the linear and probabilistic readers remain in \(\mathcal J_{0,N}^{(m)}(h_0)\). By compatibility, the limit produces \eqref{hmtfg:base:eq:continuo-discreto} without disaggregation.
\end{proof}

Subsequent recognition of the single continuum constructed through its five spectral, solenoidal, gauge, cohomological and determinantal realizations belongs to the subsequent conventional language. It does not select its states or take part in its generation.
\fi

\input{\HMTFGRoot/latex/antecedentes_comunes/continuo_cinco_construcciones_tipado.tex}
